\documentclass[11pt,a4paper]{extarticle}

\usepackage[T2A]{fontenc}
\usepackage[utf8]{inputenc}
\usepackage[russian]{babel}
\usepackage{indentfirst}
\usepackage{geometry}
\usepackage{setspace}

\geometry{left=20mm, right=10mm, top=15mm, bottom=15mm}
\onehalfspacing
\setlength{\parindent}{1.25cm}
\setlength{\parskip}{0pt}
\frenchspacing

\begin{document}
\pagestyle{empty}

\begin{center}
  \textbf{ОТЗЫВ}\\
  \textbf{научного руководителя на выпускную квалификационную работу}\\
  Киселева Никиты Сергеевича\\
  <<Методы определения достаточного размера выборки\\ в параметрических моделях>>
\end{center}

\vspace{0.3em}

Выпускная квалификационная работа Н.\,С.~Киселева связывает локальную геометрию оптимизационной поверхности параметрических моделей с задачей определения достаточного объема обучающей выборки. Тема актуальна: эмпирические закономерности масштабирования глубоких моделей не имеют согласованной теоретической интерпретации, а классические критерии стабилизации ландшафта непригодны для современных архитектур. Предложенный автором подход, в котором пробная мера согласована с подпространством наибольшей кривизны матрицы Гессе, преодолевает это ограничение и доводит количественный анализ до моделей порядка~${10^8}$ параметров. Достоверность утверждений обеспечена явно сформулированными предположениями, строгостью доказательств и воспроизводимой экспериментальной проверкой, количественно подтвердившей показатели степенного закона.

Теоретическая значимость работы состоит в построении унифицированной схемы критериев стабилизации с функцией предпочтения точек в пространстве параметров и в доказательстве для подпространственного критерия~${\Delta_2^{(D)}}$ верхней оценки~${\mathcal O(k^{-2})}$, в которой зависимость от полной размерности заменена на зависимость от размерности подпространства. Практическая значимость определяется масштабируемой алгоритмической схемой оценивания критериев на основе предлагаемого численного метода; эмпирически показано ускорение примерно в~${10^4}$ раз по сравнению с прямой оценкой методом Монте-Карло без потери точности в локальном режиме применимости.

За время совместной работы Н.\,С.~Киселев проявил себя как самостоятельный, инициативный и научно зрелый исследователь: им предложены постановка задачи, ключевые определения и теоретические результаты, спланированы и реализованы все вычислительные эксперименты, в том числе на трансформерной модели~${\sim 10^8}$ параметров, потребовавшей отдельной программной инженерии. Работа отличается академической аккуратностью: автор явно оговаривает идеализирующие предположения и не выходит за пределы того, что подтверждено доказательствами или экспериментами. Результаты прошли апробацию на ведущих российских и международных научных мероприятиях и опубликованы в рецензируемых изданиях. В качестве замечания можно отметить, что ряд содержательных направлений "--- в частности, режимы с медленным дрейфом главных собственных направлений и связь полученных оценок с эмпирическими законами масштабирования "--- сознательно оставлен за пределами работы и является естественным продолжением исследований в аспирантуре.

Считаю, что выпускная квалификационная работа Н.\,С.~Киселева в полной мере соответствует требованиям к ВКР магистров в МФТИ, заслуживает оценки~\textbf{<<отлично>>}, а её автор "--- присвоения квалификации магистра по направлению 09.04.01 <<Информатика и вычислительная техника>> и рекомендуется к продолжению научной деятельности в аспирантуре.

\vspace{0.5em}

\begin{flushleft}
\hfill\begin{tabular}{l}
\textbf{Научный руководитель}:\\
Грабовой Андрей Валериевич, к.ф.-м.н.,\\
старший научный сотрудник Лаборатории №42\\
Института проблем управления РАН им. Трапезникова,\\
доцент кафедры Интеллектуальных систем МФТИ (НИУ)
\end{tabular}
\end{flushleft}

\vfill

\begin{flushright}
    <<\underline{\hspace{1cm}}>>\ \underline{\hspace{2.5cm}}\ 2026~г.
\end{flushright}

\end{document}
